\section*{How to use this companion}
\addcontentsline{toc}{section}{How to use this companion}
This companion and the article form a single proof package. The article gives
the objects, main arguments and conclusions; the entries below identify the
calculation behind each excursion and the point at which the main argument
resumes. Unqualified numbers refer to the article; lettered numbers refer
to this companion.

\begin{center}\small
\renewcommand{\arraystretch}{1.3}
\begin{tabularx}{\textwidth}{@{}>{\raggedright\arraybackslash}p{.24\textwidth}>{\raggedright\arraybackslash}X>{\raggedright\arraybackslash}p{.25\textwidth}@{}}
\toprule
Entry from the article & Calculation here & Return to the article\\\midrule
\Cref{lem:runge} & Coefficients, integrality and the growing-degree remainder in \cref{supp:lem:runge-calculation}. & Weighted defects, \cref{prop:defects}.\\
\Cref{macro:lem:pell,lem:split-bounded} & Uniform ideal and unit counts, then localization to two square classes in A.2--A.3. & Bounded components and the deep-core estimates.\\
\Cref{prop:cutoffs,thm:scalar-rates} & Weighted prime sums, degree bounds and scalar retention in B.1--B.3. & Scalar laws or the fields of Section~5.\\
\Cref{thm:word-field,thm:typical-dictionaries} & Marginal cap, selection probabilities and collision bounds in C.1. & Uniform, typical and affine dictionaries.\\
\Cref{thm:moving-comparison}(ii) & The two Stein bounds, Poisson filling and the single tail event in C.3--C.4. & The staircase; information-adapted labels in C.5 and \cref{prop:information-labelled}.\\
\Cref{thm:resolved-kernel} & Poisson atom estimates and the resolution cost in D.1--D.2. & Conditional laws in Section~6.\\
\Cref{thm:boundary} & Prime incidences, quotient rank and the post-quadratic bound in E.1--E.5. & The microscopic mass and the prefix law.\\
\Cref{thm:relative-bulk,thm:two-clock} & The normalized error table, finite-cylinder clock and affine conditioning in E.6--E.7. & The two-source comparison and its affine variant.\\\bottomrule
\end{tabularx}
\end{center}

\paragraph{The prefix coupling and its optional complement.}
\Cref{thm:micro-run-tv} enters Appendix F. The signed input register
\cref{supp:pivot:lem:inputs} points to the one-window and two-window
estimates, local signed ranks, deep-start bound and independently proved
scalar tail. The argument then proceeds through the finite Palm--Stein
inequality, uniform Rankin estimate, reciprocal pivots and exact forcing.
It returns to the full same-grid field and
\cref{cor:micro-readout}. The restriction to the dyadic field in
\cref{cor:dyadic-micro} needs only the stability of the hard scales
under $M=4N$ and projection of integer labels. No result of Appendix G
is used along this route.
Appendix G separately proves regular-cloud statements, finite signed
Palm identities and the absolute-cumulant obstruction. Its vanishing
Palm deficit is downstream of Appendix F.
The empirical consequence \cref{cor:empirical-poisson} is proved in the
article, immediately after \cref{cor:micro-readout}: it uses concentration
of overlapping-window frequencies in the independent target and the
summable microscopic error, with no additional arithmetic input.
\Cref{rem:empirical-field-obstruction} then gives a finite-support
obstruction to the analogous local-vector comparison after fixing $f$.

\paragraph{Support and target conventions.}
For interior starts, the following probabilities hold on good supports
conditionally on the exposed prime signs.
\begin{center}\small
\begin{tabularx}{\textwidth}{@{}Xccc@{}}\toprule
Event & Values & Independent equations & Target probability\\\midrule
Run of length at least $L$ & $L+1$ & $L$ & $2^{-L}$\\
Absolute word of length $B$ & $B$ & $B$ & $2^{-B}$\\
Exact run, excess $e$ & $L+e+2$ & $L+e+1$ & $2^{-L-e-1}$\\
Exact run, excess $e$, sign $s$ & $L+e+2$ & $L+e+2$ & $2^{-L-e-2}$\\\bottomrule
\end{tabularx}
\end{center}
Summing the last row over signs gives $2^{-L-e-1}$, and then summing
over $e\ge0$ gives $2^{-L}$. With $e\le E$, the maximal support has
$Q+1$ values, $Q=L+E+1$. This is distinct from the base support
$B=L+1$ and from the exponential profile parameter $Q_B=2^B$.

\paragraph{Two different summations.}
For the graph-based spatially labelled comparison of Section~5, marks
are summed without weights before the pair estimate, leaving the
relation profile at the base length.
For an aggregated vector, the directional Stein weights are inserted
first, and the maximal-length profile is required. Section~C.3 keeps
this distinction explicit before any asymptotic absorption. Its signed
local factor can be four, not two. Full marks of base-contained starts
are not right-censored run lengths.

In the selected-pivot comparison of Appendix F the signed low types are
instead controlled on maximal supports and summed with their exact
weights. Its dimension-free bound has no inverse minimum type
probability. This is a third ledger, not a use of an aggregated
Stein factor on an unrecorded spatial field.
